% Supplementary table \label{tab:patient_level}
\begin{table*}[p]
\centering
\begin{threeparttable}

\caption{Patient-level performance of ThyroidXAgent on ZJH-8K and ThyroidXL. For classification, the patient score was the mean of the image-level malignancy probabilities and the patient label was malignant if any image of the patient was malignant. For segmentation, image-level Dice and HD95 were averaged within each patient. 95\% confidence intervals were estimated from 1{,}000 patient-level bootstrap resamples. ZJH-8K: 6{,}807 images (1{,}068 patients).}
\label{tab:patient_level}

\footnotesize
\setlength{\tabcolsep}{5.5pt}
\renewcommand{\arraystretch}{1.08}

\begin{tabular}{lccccc}
\toprule
\textbf{Cohort} & \textbf{Patients} & \textbf{AUROC (95\% CI)} & \textbf{AUPRC (95\% CI)} & \textbf{Dice, \% (95\% CI)} & \textbf{HD95, mm (95\% CI)} \\
\midrule
ZJH-8K & 1,068 & 0.9513 (0.9393--0.9633) & 0.9622 (0.9512--0.9732) & 87.91 (87.51--88.31) & 10.35 (9.96--10.74) \\
ThyroidXL & 739 & 0.9803 (0.9720--0.9886) & 0.9655 (0.9505--0.9805) & 88.40 (88.00--88.80) & 5.23 (4.73--5.73) \\
\bottomrule
\end{tabular}
\end{threeparttable}
\end{table*}
